\section{Introduction}
\label{sec:introduction}

% --- Biodiversity loss and habitat fragmentation in Switzerland ---

Biodiversity is declining at an unprecedented rate worldwide, driven by habitat 
loss, fragmentation, climate change, and land-use intensification \citep{IPBES2019}. 
Switzerland, despite its relatively small area, harbours exceptional plant 
diversity owing to its pronounced altitudinal gradients and mosaic of biogeographic 
regions. Yet even here, a third of all assessed species and half of all habitat 
types are threatened \citep{FOEN2022}, with 28\% of vascular plants considered 
endangered or regionally extinct and an additional 14\% near-threatened 
\citep{InfoFlora2024}. The main drivers---soil sealing, landscape fragmentation, 
intensive land use, and nitrogen inputs---continue to erode habitat quality and 
connectivity \citep{FOEN2022,FOEN2025}. Effective conservation planning requires 
spatially explicit predictions of where species can persist, and crucially, which 
landscape elements link suitable habitat patches into functional networks.

% --- Traditional SDMs: strengths and limitations ---

Species distribution models (SDMs) have become a cornerstone of conservation 
biogeography, relating georeferenced occurrence records to environmental 
predictors to map habitat suitability across space \citep{Elith2009,Guisan2005}. 
Classical approaches, such as generalised linear models, maximum entropy \citep{Phillips2006}, 
and ensemble methods such as Random Forests \citep{Breiman2001}, treat each 
occurrence as an independent sample drawn from an environmental niche. While powerful, 
this point-based paradigm ignores the spatial structure of landscapes: two 
patches with identical environmental conditions may differ dramatically in 
ecological value depending on their size, shape, and connectivity to neighbouring 
habitat. For species whose long-term persistence depends on seed dispersal between 
habitat patches, such as the Cypripedium calceolus, this omission can lead to 
misleading suitability estimates and suboptimal conservation recommendations.

% --- GNN-SDM framework: how graph structure captures landscape interactions ---

Graph Neural Networks (GNNs) offer a principled way to incorporate landscape 
structure into SDMs. By representing the landscape as a graph, where nodes are 
ecologically homogeneous patches and edges encode spatial adjacency. A GNN can 
learn how environmental conditions in neighbouring patches jointly influence 
local habitat suitability. \citet{Wu2025} introduced the SOM-GNN-SDM framework, 
demonstrating that a GraphSAGE architecture \citep{Hamilton2017} trained on 
landscape graphs substantially outperforms traditional SDMs for virtual and real
species across continental extents. Their approach constructs landscape patches 
via Self-Organising Maps (SOMs) \citep{Kohonen1982}, builds an adjacency graph, 
and propagates information across the graph to produce suitability predictions 
that are spatially coherent and connectivity-aware.

% --- Conservation application: graph structure for gap identification ---

Beyond improving predictive accuracy, the graph representation opens a direct 
path to conservation applications. The Swiss federal government has identified 
the lack of an ecological infrastructure that protects and connects core 
biodiversity areas as a key deficit, setting a target of 17\% of national 
territory as connected core areas by 2030, up from 13.4\% today 
\citep{FOEN2022,FOEN2025}. Because the landscape graph explicitly encodes 
connectivity, it can be analysed using circuit theory and centrality metrics to 
identify critical corridors: patches whose protection would maintain or restore 
habitat connectivity for threatened species. This transforms the SDM output from 
a static suitability map into an actionable tool for spatial conservation 
prioritisation and gap analysis, directly informing where the ecological 
infrastructure should be expanded.

% --- Our contribution ---

In this work, we adapt the GNN-SDM framework to Swiss vascular flora at 30\,m 
resolution, presenting a substantially finer grain than the original continental-scale 
application. We construct a landscape graph of 166,464 patches from 12 environmental 
features using SOM-based clustering and connected-component labelling, train 
per-species GraphSAGE models for 3,751 qualifying plant species, and benchmark 
against a Random Forest baseline. We further apply circuit-theory connectivity 
analysis to the resulting suitability maps to identify conservation priority 
areas and quantify protection gaps relative to the existing protected-area network.

% --- Research questions ---

Specifically, we address the following research questions:
\begin{enumerate}[nosep]
  \item Can GNN-SDM outperform a traditional Random Forest at patch-level species distribution modelling for Swiss plants?
  \item Does the graph structure confer greater benefit for rare and vulnerable species, whose distributions are more fragmented?
  \item What GNN architecture is best suited for flora---given that plants, unlike the mobile fauna studied by \citet{Wu2025}, have limited active dispersal?
  \item Can the graph structure be leveraged to identify spatial gaps in the protected-area network where additional conservation measures would be most beneficial?
\end{enumerate}

